%------------------------------------------------------------------------------%
\begin{question}
  Determine a posição relativa entre a reta e o plano \\
  e o ponto de interseção, se ele existir
  \[
    r:
    \begin{cases}
      x = 1 +  t \\
      y = 2 -  t \\
      z = 3 + 2t
    \end{cases}
  \]
  \[
    \gamma: x + y + z - 6 = 0
  \]
\end{question}

%------------------------------------------------------------------------------%
\begin{resolution}{Posição Relativa}
  \onslide<+->{O vetor diretor da reta é \qquad
  \(
    \vec{v} = (1, -1, 2)
  \)}

  \onslide<+->{O vetor normal do plano é \quad
  \(
    \vec{n} = (1, 1, 1)
  \)}
  \begin{align*}
    \onslide<+->{\vec{v} \cdot \vec{n}}
    \onslide<+->{& = 1 \times 1 + (- 1)1 + 2 \times 1} \\
    \onslide<+->{& = 1 - 1 + 2} \\
    \onslide<+->{& = 2} \\
    \onslide<+->{& \neq 0}
  \end{align*}
  \onslide<+->{a reta é secante ao plano}
\end{resolution}

%------------------------------------------------------------------------------%
\begin{resolution}{Ponto de Interseção}
\begin{columns}
\column{0.35\textwidth}
  \onslide<+->{Substituir a equação da reta
  \[
    \begin{cases}
      x = 1 +  t \\
      y = 2 -  t \\
      z = 3 + 2t
    \end{cases}
  \]}
  \onslide<+->{na equação do plano}
\column{0.5\textwidth}
  \begin{align*}
    \onslide<+->{x + y + z - 6              & = 0 \\}
    \onslide<+->{(1+t) + (2-t) + (3+2t) - 6 & = 0 \\}
    \onslide<+->{t - t + 2t + 1 + 2 + 3 - 6 & = 0 \\}
    \onslide<+->{2t                         & = 0 \\}
    \onslide<+->{t                          & = 0}
  \end{align*}
\end{columns}
\end{resolution}

%------------------------------------------------------------------------------%
\begin{resolution}{Ponto de Interseção}
  \onslide<+->{Coordenadas do ponto com $t=0$}
  \begin{align*}
    \onslide<+->{x & = 1 + t }
    \onslide<+->{  \;=\; 1 \\}
    \onslide<+->{y & = 2 - t}
    \onslide<+->{  \;=\; 2 \\}
    \onslide<+->{z & = 3 + 2t}
    \onslide<+->{  \;=\; 3}
  \end{align*}

  \onslide<+->{\[
    P = (1, 2, 3)
  \]}
\end{resolution}

%------------------------------------------------------------------------------%
\endinput
